\documentclass[11pt,a4paper]{report}
\usepackage[a4paper,margin=2.2cm,top=2.4cm,bottom=2.4cm]{geometry}
\usepackage{amsmath,amssymb,amsthm}
\usepackage{graphicx}

\theoremstyle{remark}
\newtheorem*{remark}{Remark}

\begin{document}

\begin{center}
Brief Description of Reasoning

Xidai Lin
\end{center}

\section*{Problem 1}
\subsection*{a}
By iterating over each element of the vector, we take the largest absolute difference as $\|\vec{x}-\tilde{\vec{x}}\|_\infty$. After computing $A\vec{x}$, a similar iteration gives $\|A\vec{x}-\vec{b}\|_\infty$.

The result is
\begin{align*}
  &\|\vec{x}-\tilde{\vec{x}}\|_\infty=0.500000 \\
&\|A\vec{x}-\vec{b}\|_\infty=0.300000  
\end{align*}



The detailed code is in \texttt{code\_of\_problem\_1a.m}.
\subsection*{b}
Similarly, substituting $\tilde{\vec{x}}=(0.33,-7.9,5.8)^\top$ gives

The result is
\begin{align*}
  &\|\vec{x}-\tilde{\vec{x}}\|_\infty=0.900000 \\
&\|A\vec{x}-\vec{b}\|_\infty=0.270000
\end{align*}

The detailed code is in \texttt{code\_of\_problem\_1b.m}.
\section*{Problem 2}
%what if A is in the field of C
\begin{proof}
First, we prove \(\rho(A^\top A)=\rho(A)^2\).

Since \(A\) is symmetric, by the Spectral Theorem there exists an
\emph{orthogonal} matrix \(Q\in M_n(\mathbb{R})\), i.e.\ \(Q^\top Q=QQ^\top=I\),
and a diagonal matrix \(\Lambda=\operatorname{diag}(\lambda_1,\dots,\lambda_n)\)
with \(\lambda_i\in\mathbb{R}\), such that
\[
A=Q^\top \Lambda Q.
\]
Hence
\[
A^\top A
=(Q^\top \Lambda Q)^\top (Q^\top \Lambda Q)
=Q^\top \Lambda Q Q^\top \Lambda Q
=Q^\top \Lambda^2 Q,
\]
where we used \(QQ^\top=I\). Thus \(A^\top A\) is \emph{orthogonally similar}
to \(\Lambda^2\), so they have the same eigenvalues:
\[
\sigma(A^\top A)=\{\lambda_1^2,\dots,\lambda_n^2\}.
\]
Therefore
\[
\rho(A^\top A)
=\max_i \lambda_i^2
=\Bigl(\max_i |\lambda_i|\Bigr)^2
=\rho(A)^2.
\]

Second, we prove that \(\rho(A^\top A)I-A^\top A\) is positive semidefinite.

Since \(Q\) is orthogonal,
\[
\rho(A^\top A)I-A^\top A
=Q^\top\bigl(\rho(A^\top A)I-\Lambda^2\bigr)Q,
\]
which is an orthogonal similarity. The matrix
\(\rho(A^\top A)I-\Lambda^2\) is diagonal with diagonal entries
\[
\rho(A^\top A)-\lambda_i^2\ge 0,\qquad i=1,\dots,n,
\]
because \(\rho(A^\top A)=\max_i\lambda_i^2\). Hence
\(\rho(A^\top A)I-\Lambda^2\) is positive semidefinite, and so is
\(\rho(A^\top A)I-A^\top A\), since orthogonal similarity preserves eigenvalues.

Finally, we prove \(\|A\|_2=\rho(A)\).

For any \(x\in\mathbb{R}^n\), the positive semidefiniteness above gives
\[
x^\top\bigl(\rho(A^\top A)I-A^\top A\bigr)x\ge 0,
\]
i.e.
\[
\|Ax\|_2^2=x^\top A^\top A x\le \rho(A^\top A)\,x^\top x
=\rho(A)^2\|x\|_2^2.
\]
Taking the supremum over \(x\neq 0\),
\[
\|A\|_2^2
=\sup_{x\neq 0}\frac{\|Ax\|_2^2}{\|x\|_2^2}
\le \rho(A)^2,
\]
so
\[
\|A\|_2\le \rho(A).
\]

Conversely, let \(\lambda\) be an eigenvalue of \(A\) with
\(|\lambda|=\rho(A)\), and let \(x\neq 0\) be a corresponding eigenvector,
i.e.\ \(Ax=\lambda x\). Then
\[
\frac{\|Ax\|_2}{\|x\|_2}
=\frac{|\lambda|\,\|x\|_2}{\|x\|_2}
=|\lambda|
=\rho(A).
\]
Since \(\|A\|_2\) is the supremum of \(\|Ax\|_2/\|x\|_2\) over all \(x\neq 0\),
we obtain
\[
\|A\|_2\ge \rho(A).
\]

Combining both inequalities,
\[
\|A\|_2=\rho(A).
\]
\end{proof}

\begin{remark}[The complex case]
For a \emph{normal} complex matrix, i.e.\ \(A^*A=AA^*\), the identity
\(\|A\|_2=\rho(A)\) still holds. By the spectral theorem for normal matrices
there exist a unitary matrix \(U\), \(U^*U=UU^*=I\), and a diagonal matrix
\(\Lambda=\operatorname{diag}(\lambda_1,\dots,\lambda_n)\) such that
\[
A=U\Lambda U^*.
\]
Then
\[
A^*A=(U\Lambda U^*)^*(U\Lambda U^*)=U\Lambda^*\Lambda U^*
=U\operatorname{diag}(|\lambda_1|^2,\dots,|\lambda_n|^2)U^*,
\]
so the eigenvalues of \(A^*A\) are \(|\lambda_1|^2,\dots,|\lambda_n|^2\), and
therefore
\[
\|A\|_2=\sqrt{\rho(A^*A)}=\max_i|\lambda_i|=\rho(A).
\]
In particular, every Hermitian matrix \(A^*=A\) is normal, so this extends the
real symmetric case above. By contrast, a complex \emph{symmetric} matrix
\(A^\top=A\) is in general neither Hermitian nor normal, so
\(\|A\|_2=\rho(A)\) need not hold for it.
\end{remark}

\section*{Problem 3}
The general code is in \texttt{gauss.m}, which implements Gaussian elimination with scaled partial pivoting.
\subsection*{a}
Interchange E1 and E2, then implement Gaussian elimination.
The final result is:
\begin{align*}
     x1 &= 10.0000 \\
     x2 &= 1.0000
\end{align*}
The detailed code is in \texttt{code\_of\_problem3a.m}.
\subsection*{b}
Similarly, the final answer is
\begin{align*}
  x1 &=0.0000 \\
  x2 &=10.0000 \\
  x3 &= 0.1429
\end{align*}
\section*{Problem 4}
The code is in \texttt{jacobi.m}.
\subsection*{a}
The result is

The 1st iteration \\
    1.2500 \\
   -1.3333 \\
    0.2000 \\

The 2nd iteration \\
  1.6333 \\
   -0.9833 \\
    0.2333 \\

The 3rd iteration \\
    1.5542 \\ 
   -0.8667 \\
   -0.0600 \\
\subsection*{b}
The result is

The 1st iteration \\
    -2 \\
     2 \\
     0 \\

The 2nd iteration \\
    -1 \\
     1 \\
    -1 \\

The 3rd iteration \\
   -1.7500 \\
    1.7500 \\
   -0.5000 \\
\section*{Problem 5}

The detailed code is in \texttt{seidel.m} and \texttt{Jacobi\_method.m}.
\subsection*{a}
Gauss--Seidel algorithm: \\
The 6th iteration \\
    0.0354 \\
   -0.2368 \\
    0.6578 \\
Jacobi algorithm: \\
The 10th iteration \\
    0.0351 \\
   -0.2369 \\
    0.6578 \\
\subsection*{b}
Gauss--Seidel algorithm: \\
The 4th iteration \\
    0.9957\\
    0.9579\\
    0.7916\\
Jacobi algorithm: \\
The 7th iteration\\
    0.9958\\
    0.9579\\
    0.7916\\
\end{document}
